\Needspace{8\baselineskip}
\SourceLeaf{94}%
\noindent Die Definitionsgleichung für $\lambda'$ ergibt nach kurzer Umrechnung\hfill\\
zunächst \mbox{$\bigl(\frac{\lambda'}{m'}\bigr)^{2}=c^{2}$}. Hier sind

\[
\frac{\lambda'}{m'}=-c,\quad\text{jedoch weiter}\quad
\frac{\lambda''}{m''}=c',\qquad \frac{\lambda'''}{m'''}=c'',\ \ldots
\]

\noindent zu setzen, um die Formeln der vorliegenden Notiz zu erhalten. Folglich gelten die
Gleichungen.

\[
\nu=-u,\qquad P=\frac{(1+t)^{2}}{4t}\,\frac{1-u}{a^{(\infty)}},\qquad
Q=\frac{(1-t)^{2}}{4t}\,\frac{1+u}{a^{(\infty)}},
\]

\noindent wie aus den Schlussformeln in Art.\,19 in\,der\,genannten Abhandlung\ hervorgeht. Durch Addition und
Subtraktion dieser letzten beiden Formeln erhält man

\[
A^{(0)}=\frac{h-u}{a^{(\infty)}},\qquad A^{(1)}=\frac{1-hu}{a^{(\infty)}}.
\]

Zweitens fällt die hier gelöste Aufgabe,
$(h-\cos\varphi)^{-\frac{3}{2}}$ zu entwickeln, unter das allgemeine Problem, das in der Dissertatiuncula (S.\,35 ff.)
behandelt wird. Unter Nr.\,7 werden dort die dreigliedrigen Beziehungen von \textsc{Euler}
aufgestellt, mit denen man aus den beiden ersten Koeffizienten
$A^{(0)}$, $A^{(1)}$ nacheinander alle folgenden $A^{(2)}$,
$A^{(3)}$, $\ldots$\ berechnen kann. Für die Quotienten $B^{(k)}$ der
$A$ ergeben sich aus den Rekursionsformeln von \textsc{Euler} die hier von \textsc{Gauss} aufgestellten
Formeln für den Zusammenhang zwischen $B^{(k)}$ und $B^{(k+1)}$.

\Needspace{20\baselineskip}
Schließlich ist noch erwähnenswert, dass der für $B^{(k)}$ angegebene
Kettenbruch unmittelbar unter die allgemeine Formel [25] in Art.\,12 der »Disquisitiones
generales circa seriem infinitam etc.« (Gesammelte Werke III S.\,134) fällt, in der
der Quotient zweier hypergeometrischer Funktionen in einen Kettenbruch entwickelt wird.
In der dort erläuterten Notation würde die Formel für $B^{(k)}$
so lauten.

\[
B^{(k)}=\frac{k+\frac{3}{2}}{k+1}\,t\mathord{.} G\bigl(-\tfrac{1}{2},k-\tfrac{1}{2},k+1;t^{2}\bigr)
\]

\noindent oder auch:

\[
B^{(k)}=\frac{k+\frac{3}{2}}{k+1}\,t\mathord{.}
\frac{F\bigl(-\tfrac{1}{2},k+\tfrac{1}{2},k+2;t^{2}\bigr)}{F\bigl(-\tfrac{1}{2},k-\tfrac{1}{2},k+1;t^{2}\bigr)}.
\]

\begin{flushright}
\textsc{Fricke.}
\end{flushright}
